\scanpage{32}
\mspage{6 - 40 -}\unsure{,}\hfill n\textsuperscript{o} 392

c) Si $\underline e(E)$ est un idéal principal $A\alpha$, montrer que $E$ contient un facteur direct isomorphe à $A/A\alpha$, (écrire $E$ comme quotient de $A^r$ par un sous-module $R$ et noter qu'il y a dans $R$ un élément de la forme $\alpha y$, o\rlap{\handwritten{\`{\phantom{u}}}}u $y$ est un élément d'une base de $A^r$).

d) Soit $\underline d'_h(E)$ l'annulateur de $\Lambda^{h+1}E$ pour $0\leq h\leq r-1$. Montrer que les conditions suivantes sont équivalentes :

$\alpha$) Les idéaux $\underline d_h(E)$, $(0\leq h\leq r-1)$, sont principaux.

$\beta$) Les idéaux $\underline d'_h(E)$, $(0\leq h\leq r-1)$, sont principaux.

$\gamma$) $E$ est somme directe de $r$ modules isomorphes à $A/A\lambda_h$, o\rlap{\handwritten{\`{\phantom{u}}}}u $\lambda_{h+1}$ divise $\lambda_h$ pour $0\leq h\leq r-1$.

En outre, lorsque ces conditions sont vérifiées, on a $\underline d'_h(E)=A\lambda_h$ pour $0\leq h\leq r-1$, et $\underline d_h(E)=\lambda_h\lambda_{h+1}\ldots\lambda_{r-1}A$ pour $0\leq h\leq r-1$. (Procéder par récurrence sur $r$, en utilisant c)).

e) On suppose en outre que $A$ est intègre. Montrer que si $\underline p$ est un idéal premier de $A$ tel que $E/\underline p E$ soit un $(A/\underline p)$-module libre et $E_p$ un $A_{\underline p}$-module libre, alors $E$ est un $A$-module libre, (en utilisant b), montrer que $\gamma(E_{\underline p})=\gamma(E)$ et que $\underline e(E_{\underline p})=\underline e(E)$).

23) Soient $A$ un anneau, $E$ un $A$-module de type fini, $r$ le plus petit des entiers $k$ tels que $\Lambda^{r+1}E=0$, $\underline e(E)$ l'annulateur de $\Lambda^r E$. Montrer que pour tout idéal $\underline a\supset\underline e(E)$ dans $A$, le $(A/\underline a)$-module $E/\underline a E$ est plat, (se ramener au cas o\rlap{\handwritten{\`{\phantom{u}}}}u $A$ est un anneau local, et utiliser l'exerc. 22 b)). En déduire que si $\underline r$ est le radical de $A$ et si, pour tout idéal maximal $\underline m$ de $A$, la dimension du $(A/\underline m)$-espace vectoriel $E/\underline m E$ est la même, alors $E/\underline r E$ est un $(A/\underline r)$-module plat.

\begin{center}
-:-:-:-
\end{center}
% archive sheet number: 31

